


% 除非校方修改了論文格式 (margins, header, footer, 浮水印)
% 或者需要增加所用的 LaTeX 套件，
% 或者要改預設中文字型、編碼
% 否則毋須修改本檔內容
% 論文撰寫，請修改以 my_  開頭檔名的各檔案




\usepackage[numbers]{natbib} % for smart citation
\usepackage{geometry}  % for easy margin settings
\usepackage[titletoc,title]{appendix} % for appendix formatting
%\usepackage{subfigure}  % for subfigure  因為有衝突
%\usepackage[dvipdfm]{graphicx}  % for graphic   using eps
%\usepackage[xetex]{graphicx}
\usepackage{graphicx}  % for graphic   using eps

%\usepackage{epstopdf} % 當使用pdflatex時打開，如使用latex則不需開啟，此功能為將xxx.eps 自動判讀為XXX.pdf


%\usepackage{algorithmic}  %演算法使用
\usepackage{algorithm}
\usepackage{algorithmicx}
\usepackage{algpseudocode}

%
% margins setting
\geometry{verbose,a4paper,tmargin=2.3cm,bmargin=3.5cm,lmargin=2.5cm,rmargin=3cm}
%
\usepackage{amsmath} % 各式 AMS 數學功能
%\usepackage{amssymb} % 各式 AMS 數學符號 (如果amssymb和Xunicode產生衝突，請將amssymb移到最前面)
\usepackage{mathrsfs} %草寫體數學符號，在數學模式裡用 \mathscr{E} 得草寫 E
\usepackage{listings} % 程式列表套件
%
% listing setting
\lstset{breaklines=true,% 過長的程式行可斷行
extendedchars=false,% 中文處理不需要 extendedchars
texcl=true,% 中文註解需要有 TeX 處理過的 comment line, 所以設成 true
comment=[l]\%\%,% 以雙「百分號」做為程式中文註解的起頭標記，配合 MATLAB
basicstyle=\small,% 小號字體, 約 10 pt 大小
commentstyle=\upshape,% 預設是斜體字，會影響註解裏的英文，改用正體
%language=Octave % 會將一些 octave 指令以粗體顯示
}

\usepackage{url} % 在文稿中引用網址，可以用 \url{http://www.ntust.edu.tw} 方式

% 插圖套件 graphicx
% 使用者工作流程是用 pdftex 還是 latex + dvipdfmx?
% 視情況而有不同的參數
% 這裡作自動判斷
% 參考自
% http://www.tex.ac.uk/cgi-bin/texfaq2html?label=ifpdf
%\newcommand\mydvipdfmxflow{dvipdfmx}
%\newcommand\mypdftexflow{pdftex}
%
%\ifx\pdfoutput\undefined
%  % not running pdftex
%  \usepackage[dvipdfm]{graphicx}
%  \newcommand\myworkflow{dvipdfmx}  % set the flag for hyperref
%\else
%  \ifx\pdfoutput\relax
%    % not running pdftex
%    \usepackage[dvipdfm]{graphicx}
%    \newcommand\myworkflow{dvipdfmx}  % set the flag
%  \else
%    % running pdftex, with...
%    \ifnum\pdfoutput>0
%      % ... PDF output
%      \usepackage[pdftex]{graphicx}
%      \newcommand\myworkflow{pdftex}  % set the flag
%    \else
%      %...DVI output
%      \usepackage[dvipdfm]{graphicx}
%      \newcommand\myworkflow{dvipdfmx}  % set the flag
%    \fi
%  \fi
%\fi

\usepackage{fancyhdr}  % 借用增強功能型 header 套件來擺放浮水印 
% (佔用了 central header)
% 不需要浮水印的使用者仍可利用此套件，產生所需的 header, footer
%
% 啟動 fancy header/footer 套件
\pagestyle{fancy}
\fancyhead{}  % reset left, central, right header to empty
\fancyfoot[C]{\thepage} %中間 footer 擺放頁碼
\renewcommand{\headrulewidth}{0pt} % header 的直線; 0pt 則無線

% 如果不需要任何浮水印，則請把下列介於 >>> 與 <<< 之間
% 的文字行關掉 (行首加上百分號)
%% 浮水印 >>> 
%
% this file is encoded in utf-8
% v1.7
% 如果浮水印不是全篇需要，請把下列介於 >>> 與 <<<
% 的「全篇浮水印專用碼」關掉 (行首加百分號)
% 參考自 Keith Reckdahl 寫的 "Using Imported Graphics in LATEX2e" (epslatex.pdf) p.39
% 如果只有特定頁需要浮水印
% 則依該頁屬性使用下列之一的命令 
% 普通頁命令 \thispagestyle{WaterMarkPage}
% plain 頁命令 \thispagestyle{PlainWaterMarkPage}
% empty 頁命令 \thispagestyle{EmptyWaterMarkPage}


% 將重複使用的浮水印章
% 圖檔是 my_watermark.xxx
% 副檔名可以不加，可以是 latex 系統能處裡的任何格式：pdf, gif, png, jpg, eps, ...
% 某些圖檔格式在某些工作流程可能需要作前置處裡。
% 例如，pdflatex 無法直接處理 eps 檔
%  latex + dvipdfmx 無法直接處理 pdf, gif, png, jpg, 需要用 ebb 小工具程式
%  對圖檔產生 .bb 對應檔。
% old code
%\newsavebox{\mywatermark}
%\sbox{\mywatermark}{\includegraphics[keepaspectratio,%
%height=0.8\textheight,%
%width=0.8\textwidth]{my_watermark}}
% new code
\newsavebox{\mywatermark}
\sbox{\mywatermark}{\includegraphics[keepaspectratio,
width=2.8cm]{watermark/ncku_watermark.eps}}


% 將 central header 擺放浮水印的巨集指令
\newcommand{\PlaceWaterMark}{\fancyhead[C]{\setlength{\unitlength}{1in}%
\begin{picture}(0,0)%
%\put(-2.2,-6){\usebox{\mywatermark}}% old code
%\put(-0.5,-6.2){\usebox{\mywatermark}}% new code
\put(-0.5,-5.3){\usebox{\mywatermark}}% new code
\end{picture}}%
}

\fancyhead{}  % reset left, central, right header to empty
%% 如果不需整篇論文都要浮水印
%% 則下面  >>> 與 <<< 之間的程式碼請關閉
%% >>> 全篇浮水印
\PlaceWaterMark  % 每一頁都有浮水印 (除了 plain、empty 頁以外)

% 重新定義 plain 頁面
% 每張 plain 頁面 (每一章的第一頁) 也加浮水印

\fancypagestyle{plain}{%
\fancyhead{}%
\PlaceWaterMark%
\fancyfoot{}%
\fancyfoot[C]{\thepage}
\renewcommand{\headrulewidth}{0pt}%
\renewcommand{\footrulewidth}{0pt}%
}
%% <<< 全篇浮水印

%% 如果只有一、兩頁需要有浮水印
%% 可以在該頁 (有頁碼) 使用 \thispagestyle{WaterMarkPage}
%% 此命令不影響原有的 header、footer
\fancypagestyle{WaterMarkPage}{%
\PlaceWaterMark%
}

%% 如果只有一、兩頁 plain 頁需要有浮水印 (如 摘要、自傳等)
%% 可以在該頁 (有頁碼) 使用 \thispagestyle{PlainWaterMarkPage}
%% 只有頁碼與浮水印，沒有其他的 header、footer
%% 等同於 plain page style + water mark
\fancypagestyle{PlainWaterMarkPage}{%
\fancyhead{}%
\PlaceWaterMark%
\fancyfoot{}%
\fancyfoot[C]{\thepage}
\renewcommand{\headrulewidth}{0pt}%
\renewcommand{\footrulewidth}{0pt}%
}

%% 如果只有一、兩頁 empty 頁需要有浮水印 (如封面、書名頁)
%% 可以在該頁 (無頁碼) 使用 \thispagestyle{EmptyWaterMarkPage}
%% 等同於 empty page style + water mark
\fancypagestyle{EmptyWaterMarkPage}{%
\fancyhead{}%
\PlaceWaterMark%
\fancyfoot{}%
\renewcommand{\headrulewidth}{0pt}%
\renewcommand{\footrulewidth}{0pt}%
}

%% <<< 浮水印



% global page layout
\newcommand{\mybaselinestretch}{1.5}  %行距 1.5 倍 + 20%, (約為 double space)
\renewcommand{\baselinestretch}{\mybaselinestretch}  % 論文行距預設值
\parskip=2ex  % 段落之間的間隔為兩個 x 的高度
\parindent = 24Pt  % 段首內縮由 CJK 控制，所以這裡就設成不內縮

% Chinese enviroment
\ifdefined \useChinese
    % 中文編號 / Chinese numbering
    \usepackage{zhnumber}
    %\renewcommand{\thechapter}{\zhnum{chapter}}  % 和 section 配合仍然很怪
\fi

%%%%%%%%%%%%%%%%%%%%%%%%%%%%%
%  end of preamble
%%%%%%%%%%%%%%%%%%%%%%%%%%%%%
